\section{Fundamentos de estatística e probabilidade}
\label{sec:estatistica}

Para decidir o que é ``fora do normal'' precisamos descrever o normal com
números. Estatística é a linguagem disso.

\subsection{Distribuição, média e desvio-padrão}
\begin{definicao}[Distribuição]
A distribuição de um conjunto de valores diz \emph{quais valores aparecem e com
que frequência}. Visualiza-se com um \textbf{histograma}: barras cuja altura é a
contagem de valores em cada faixa.
\end{definicao}

\begin{definicao}[Média e desvio-padrão]
A \textbf{média} $\mu = \frac{1}{n}\sum_i x_i$ é o centro de massa dos dados.
O \textbf{desvio-padrão} $\sigma = \sqrt{\frac{1}{n}\sum_i (x_i-\mu)^2}$ mede o
espalhamento típico em torno da média: pequeno $\sigma$ = valores agrupados,
grande $\sigma$ = valores dispersos.
\end{definicao}

\begin{intuicao}
As alturas de pessoas adultas têm média $\approx 1{,}70$\,m e desvio-padrão
$\approx 0{,}10$\,m. Alguém com $1{,}72$\,m é banal; com $2{,}10$\,m é
\emph{raro} --- está a quatro desvios-padrão da média. ``Quantos desvios-padrão
de distância'' é uma medida natural de ``quão incomum''.
\end{intuicao}

\subsection{Percentil: cortando pela frequência, não pela fórmula}
\begin{definicao}[Percentil]
O percentil $p$ é o valor abaixo do qual caem $p\%$ das observações. O
percentil 95 (p95) é o valor que deixa 95\% dos dados abaixo e 5\% acima.
\end{definicao}

\begin{exemplo}
Em 100 provas com notas ordenadas, o p95 é (aproximadamente) a 95\textsuperscript{a}
nota: só 5 alunos tiraram acima dela. Quem passa do p95 está no topo dos 5\%.
\end{exemplo}

\begin{intuicao}[Por que percentil e não ``média + k desvios''?]
A regra ``$\mu + k\sigma$'' pressupõe que os dados seguem o sino da curva
normal. Retornos financeiros \emph{não} seguem: têm \textbf{caudas pesadas}
(eventos extremos são mais comuns do que a curva normal preveria). O percentil
é \textbf{não paramétrico} --- não assume formato nenhum, apenas conta
frequências --- por isso é mais robusto aqui.
\end{intuicao}

\begin{naprat}
O limiar que separa ``normal'' de ``anômalo'' (Seção~\ref{sec:deteccao}) é o
\textbf{percentil 95 do erro de reconstrução}: marcamos como suspeitos os 5\%
de maior erro. Isso embute uma \emph{taxa-base} de $\approx$5\% de marcações
mesmo em dados normais --- um detalhe que volta a importar na avaliação
(Seção~\ref{sec:avaliacao}).
\end{naprat}

\subsection{Normalização: colocar tudo na mesma escala}
Redes neurais aprendem melhor quando os números de entrada vivem em uma faixa
pequena e comparável (tipicamente $[0,1]$). A \textbf{normalização Min--Max}
faz isso linearmente:
\[
x' = \frac{x - x_{\min}}{x_{\max} - x_{\min}} .
\]
O menor valor vira 0, o maior vira 1, e o resto fica proporcional no meio.

\begin{lstlisting}[caption={Min--Max com scikit-learn.}]
from sklearn.preprocessing import MinMaxScaler
scaler = MinMaxScaler()
scaler.fit(treino.reshape(-1, 1))     # aprende min e max -- so do treino!
treino_norm = scaler.transform(treino.reshape(-1, 1))
teste_norm  = scaler.transform(teste.reshape(-1, 1))
\end{lstlisting}

\begin{intuicao}
O detalhe crucial está no comentário do código: o \texttt{scaler} aprende
$x_{\min}, x_{\max}$ \textbf{apenas no treino}. Por quê? Porque usar o mínimo e
o máximo do período de teste seria espiar o futuro. Esse é o tema da próxima
seção.
\end{intuicao}
